\section{Discussion and conclusion}
\label{sec:discussion}

This work develops a unified variational framework for the analysis of finite-pool
Chebyshev greedy algorithms (CGA) for nonlinear PDEs. The atomic source condition and
directional smoothness yield an $O(m^{-1})$ energy bound for the ideal
continuous-dictionary CGA. For the practical finite-pool algorithm, the analysis further
separates the effects of continuous-dictionary coverage, score-evaluation error, radial
stationarity, and inexact coefficient optimization, and identifies relative accuracy
conditions under which the same energy order is retained. This provides a quantitative
description of how finite-pool approximation, numerical score evaluation, and inexact
active-space solves influence the greedy descent mechanism.

For the standard pure, regularized, and reaction $p$-energies with $p\geq2$, we further
establish quantitative relations among the energy gap, the natural $V$-distance, and the
Sobolev error. In particular, an energy exponent $\alpha$ leads to the corresponding
natural-distance and $W^{1,p}$-error exponents $\alpha/2$ and $\alpha/p$, respectively.
This allows energy estimates for the greedy algorithm to be translated directly into error
measures adapted to the nonlinear geometry of the underlying PDEs.

Building on this structure, two refined analytical mechanisms are developed. The local
Hessian-transfer analysis connects the nonlinear problem with greedy approximation in a
frozen Hilbert metric, while the finite-trajectory analysis works directly with the nested
spaces generated by the computation and produces explicit projection bounds through either
an approximate source representation or normalized innovations. For the quartic energy, an
exact energy decomposition further connects these projection estimates with nonlinear
energy errors. Consequently, even on computational ranges where a uniform asymptotic
theory is not yet available, the analysis provides quantitative and directly testable
explanations of the observed algorithmic trajectories.

The one-dimensional pure and regularized $p=4$ experiments support this
finite-trajectory framework. On the calibration range $N=8$--$32$, both the
source-based and fractional-innovation bounds enclose the computed errors.
For the prescribed benchmark $\alpha=6$, the conditional comparison powers
on this range are $6$, $3$, and $3/2$ for the energy, natural distance, and
$W^{1,4}$ error, respectively. The projection and quartic identities remain
accurate under two quadrature rules on the continuation to $N=64$.
The normalized-innovation sufficient condition holds for 27 of 32 continuation
transitions in the pure model and 28 of 32 in the regularized model.
Numerical coverage by a conservative expression on the continuation does not
establish the same theorem guarantee throughout that range.

The equal-coefficient comparisons use refitted CGA prefixes and fully
archived FEM and RFM states with the same metric-specific quadrature policy.
Their computed CGA and RFM endpoints, together with interpolated FEM reference
values where indicated, describe the approximation errors for the specified problems and valid samples. They do not establish uniform superiority over
all widths, excluded realizations, or computational-work budgets. The
componentwise gradient seminorm, full Sobolev norm, and natural distance are
reported as distinct metrics rather than interchangeable relative errors.

Overall, the theoretical and numerical results show that finite-pool CGA admits a coherent
variational analysis while providing effective low-dimensional adaptive approximations for
nonlinear PDEs with $p$-structure. The present framework also identifies concrete and
testable conditions for further developments, including continuous-dictionary coverage,
entropy bounds for the restarted dictionary, and invariance of the Hessian neighborhood.
Establishing these conditions uniformly, together with reliable numerical enclosures and
comparisons under a common computational-work measure, would provide a natural route from
the present finite-range theory toward a more complete asymptotic convergence and
complexity theory.
